%%%%%%%%%%%%%%%%%%%%%%%%%%%%%%%%%%%%%%%%%%%%%%%%%%%%%%%%%%%%%%%%%%%%%
%%                                                                 %%
%% Please do not use \input{...} to include other tex files.       %%
%% Submit your LaTeX manuscript as one .tex document.              %%
%%                                                                 %%
%% All additional figures and files should be attached             %%
%% separately and not embedded in the \TeX\ document itself.       %%
%%                                                                 %%
%%%%%%%%%%%%%%%%%%%%%%%%%%%%%%%%%%%%%%%%%%%%%%%%%%%%%%%%%%%%%%%%%%%%%

%%\documentclass[referee,sn-basic]{sn-jnl}% referee option is meant for double line spacing

%%=======================================================%%
%% to print line numbers in the margin use lineno option %%
%%=======================================================%%

%%\documentclass[lineno,sn-basic]{sn-jnl}% Basic Springer Nature Reference Style/Chemistry Reference Style

%%======================================================%%
%% to compile with pdflatex/xelatex use pdflatex option %%
%%======================================================%%

%%\documentclass[pdflatex,sn-basic]{sn-jnl}% Basic Springer Nature Reference Style/Chemistry Reference Style

%%\documentclass[sn-basic]{sn-jnl}% Basic Springer Nature Reference Style/Chemistry Reference Style
%%\documentclass[sn-mathphys]{sn-jnl}% Math and Physical Sciences Reference Style
%%\documentclass[sn-aps]{sn-jnl}% American Physical Society (APS) Reference Style
%%\documentclass[sn-vancouver]{sn-jnl}% Vancouver Reference Style
%%\documentclass[sn-apa]{sn-jnl}% APA Reference Style
%%\documentclass[sn-chicago]{sn-jnl}% Chicago-based Humanities Reference Style
%%\documentclass[sn-standardnature]{sn-jnl}% Standard Nature Portfolio Reference Style
\documentclass[default]{sn-jnl}% Default
% \documentclass[default,iicol]{sn-jnl}% Default with double column layout

%%%% Standard Packages
%%<additional latex packages if required can be included here>
%%%%
\usepackage{epsfig}
\usepackage{graphicx}
\usepackage{amsmath}
\usepackage{amssymb}
\usepackage{booktabs}
\usepackage{comment}
\usepackage{color}

\usepackage{latexsym}
\usepackage{makecell,diagbox}
\usepackage{algorithm}
\usepackage{listings}

\usepackage{amsfonts}
\usepackage{colortbl}
\usepackage{microtype}
\usepackage{multirow}
\usepackage{adjustbox}
\usepackage{enumitem}
\usepackage{float}
\usepackage{bm}
\usepackage{bbm} % \mathbbm

\usepackage{bbding}
\usepackage{pifont}
\usepackage{wasysym}
\usepackage{xcolor}
% \usepackage{caption}

%%%%%=============================================================================%%%%
%%%%  Remarks: This template is provided to aid authors with the preparation
%%%%  of original research articles intended for submission to journals published 
%%%%  by Springer Nature. The guidance has been prepared in partnership with 
%%%%  production teams to conform to Springer Nature technical requirements. 
%%%%  Editorial and presentation requirements differ among journal portfolios and 
%%%%  research disciplines. You may find sections in this template are irrelevant 
%%%%  to your work and are empowered to omit any such section if allowed by the 
%%%%  journal you intend to submit to. The submission guidelines and policies 
%%%%  of the journal take precedence. A detailed User Manual is available in the 
%%%%  template package for technical guidance.
%%%%%=============================================================================%%%%

\jyear{2022}%

%% as per the requirement new theorem styles can be included as shown below
\theoremstyle{thmstyleone}%
\newtheorem{theorem}{Theorem}%  meant for continuous numbers
%%\newtheorem{theorem}{Theorem}[section]% meant for sectionwise numbers
%% optional argument [theorem] produces theorem numbering sequence instead of independent numbers for Proposition
\newtheorem{proposition}[theorem]{Proposition}% 
%%\newtheorem{proposition}{Proposition}% to get separate numbers for theorem and proposition etc.

\theoremstyle{thmstyletwo}%
\newtheorem{example}{Example}%
\newtheorem{remark}{Remark}%

\theoremstyle{thmstylethree}%
\newtheorem{definition}{Definition}%

\raggedbottom
%%\unnumbered% uncomment this for unnumbered level heads

\begin{document}

\title[]{Summary of Changes and Response Letter}

%%=============================================================%%
%% Prefix	-> \pfx{Dr}
%% GivenName	-> \fnm{Joergen W.}
%% Particle	-> \spfx{van der} -> surname prefix
%% FamilyName	-> \sur{Ploeg}
%% Suffix	-> \sfx{IV}
%% NatureName	-> \tanm{Poet Laureate} -> Title after name
%% Degrees	-> \dgr{MSc, PhD}
%% \author*[1,2]{\pfx{Dr} \fnm{Joergen W.} \spfx{van der} \sur{Ploeg} \sfx{IV} \tanm{Poet Laureate} 
%%                 \dgr{MSc, PhD}}\email{iauthor@gmail.com}
%%=============================================================%%


\maketitle

%%%%%%%%% BODY TEXT

We wish to extend our heartfelt appreciation to the editors and reviewers for their meticulous evaluation of our manuscript titled ``CLIP-Adapter: Better Vision-Language Models with Feature Adapters''. The insightful feedback are immensely valuable in enhancing the overall quality of our paper. We have carefully incorporated all the suggestions and thoroughly revised our manuscript. The following section provides a concise summary of the modifications we have made and our detailed responses to address their concerns, point by point. Accordingly, this document is structured as follows:
\begin{itemize}
    \item Section~\ref{sec:summary}: Summary of Changes

    \item Section~\ref{sec:editor}: Response to Editor
    
    \item Section~\ref{sec:reviewer_1}: Response to Reviewer \#1
    
    \item Section~\ref{sec:reviewer_2}: Response to Reviewer \#2
    
    \item Section~\ref{sec:reviewer_3}: Response to Reviewer \#3
    
    \item Annotated Revision (with highlighted text)
\end{itemize}


\section{Summary of Changes}
\label{sec:summary}
We carefully revised the paper to thoroughly address all questions and suggestions raised by the reviewers, providing point-by-point responses. Specifically, our revisions encompassed:

\begin{enumerate}
    \item We added Figure 1 to illustrate the difference between CoOp and our CLIP-Adapter, especially for how they differ in gradient back-propagation.
    \item In Section 1, we clearly stated our motivation of only adding two layers of linear layers after the text or vision backbones.
    \item We polished Section 2 to improve the descriptions of related works, especially for including more suggested related works in the field of ``natural
language prompts for vision-language pretrained models''.
    \item We provided our efficiency evaluation setup in Table 1 as well as explanations of why CLIP-Adapter is more efficient than CoOp in Section 4.2.3.
    \item We added Table 2 to provide new comparison of adapters inserted into different layers of CLIP to show that our CLIP-Adapter's design is reasonable. We also provided detailed analysis for where to insert CLIP-Adapter in Section 4.2.6.
    \item In Table 3, we offered comparison of CLIP-Adapter with houlsby-adapter~\cite{houlsby2019parameter} and MAM adapter~\cite{he2022towards}. We also gave detailed analysis in Section 4.2.7.
    \item We provided comparison of CLIP-Adapter and ELEVATER
benchmark baselines in Table 4 and corresponding discussions in Section 4.2.8.
    \item We reorganized figures and tables to improve the presentation of our revised manuscript.
\end{enumerate}


\section{Response to Editor}
\label{sec:editor}

\textbf{Point 1:} The manuscript was reviewed by three experts in the field. The reviewers expressed several concerns regarding the unclear motivation, comparisons to common adapter techniques, and baseline setup with different backbones. Considering the nature of the reviewers' comments, it is felt that such concerns can be addressed in a revision of the paper. The authors are encouraged to address all the reviewers' comments in the revision.

\noindent \textbf{Reply:} Thank you very much for your hard efforts on handling our submission. These constructive suggestions are highly helpful for us to improve our paper. We followed all the comments to revise our paper thoroughly. For further details, please refer to the following response.



\section{Response to Reviewer \#1}
\label{sec:reviewer_1}

\textbf{Point 1:} What is the motivation and justification of only adding two layers of linear layers after the text and vision backbones? Since the authors present the method as ``CLIP-adapter'', one would expect the proposed method is the most suitable adapter technique for CLIP. Although the authors has provided detailed ablation studies, the current submission does not reflect this issue (if I wasn't missing anything). More specifically, the authors are suggested to compare common adapter technique with the proposed method, for example, houlsby-adapter~\cite{houlsby2019parameter}, and MAM adapter~\cite{he2022towards}.

\noindent \textbf{Reply:} The motivation of our CLIP-Adapter is that the original pretrained encoder of CLIP has already been equipped with strong representation capabilities, so it only requires a lightweight adaption in the form of residuals. As presented in Table 2, our approach achieves superior performance with minimal computational cost when inserted at the end (i.e., the 12th layers).

Compared to other methods (houlsby-adapter and MAM adapter), as shown in Table 3, our approach outperforms them in terms of both performance and efficiency. This is because we can preserve the knowledge obtained from CLIP by using the residual-style in the end, and do not need to back-propagate the gradients through the encoders.
However, existing methods often adopt the non-residual forms and insert their adapters densely into every middle layer of CLIP's encoders, which may result in information loss and overfitting.

\vspace{5pt}
\noindent \textbf{Point 2:} The authors only investigate methods to adapt CLIP model, a particular pretrained vision-language model. The proposed method is not as generalizable to other more common image classification backbones, e.g. ViT, Swin, or ResNet.

\noindent \textbf{Reply:} It is important to note that our residual method is built upon pretrained backbones with knowledge stored in the text encoder, which allows the ability for zero-shot visual recognition. ViT, Swin, and ResNet pretrained by ImageNet contain no text encoder and are thus not suitable for our adapter method.
Since CLIP is a landmark approach in the field of vision-language understanding, therefore, we focused on adapting CLIP in our paper. We have also experimented CLIP-Adapter with different CLIP's visual backbones in Table 8, e.g., ResNet, ViT-B/32, and ViT-B/16.
% However, it would be interesting to investigate how this method could be extended to other pretrained models with text encoders in the future.

\vspace{5pt}
\noindent \textbf{Point 3:} Table 1 shows that CoOp needs 14h 40min training time, while fully finetuning CLIP backbones and CLIP-adapters only takes 3h+ (Table 10). What is the reason for such counter-intuitive results? Since CoOp only needs to propagate the text backbone.

\noindent \textbf{Reply:} In our final version of the proposed method, the adapter is added after the encoder, which allows for efficient computation and back-propagation without the need to calculate gradients for the entire encoder. In contrast, the CoOp method has learnable parameters (soft prompts) located at the front of the input to feed in the text encoder, which requires the computation and storage of gradients for the entire encoder during back-propagation, leading to increased computational overhead and memory usage. As illustrated in the Figure 1, there are differences in the computational graph and gradient calculation between the two methods. Overall, our approach enables more efficient and scalable training compared to the CoOp method.

\vspace{5pt}
\noindent \textbf{Point 4:} 
% Table 1 indicates that there is no GPU needed for linear-probing. 
Did the author use CPU only for the linear-probing experiments? Table 1 and 10 should use the same GPU(s) from the same machine for a fairer comparison.

\noindent \textbf{Reply:} Thank you for your attention to the details of our experiments. Yes, we used the Scikit-learn library and the NumPy library for linear probing on CPU, which is consistent with the method used in the CoOp paper. 

To ensure a fair comparison between the results, we used the same GPU on the same machine for both Table 1 and Table 10. Throughout our experiments, we were mindful of and addressed any potential hardware and software differences that may affect the results, to ensure their reliability and consistency. 

\vspace{5pt}
\noindent \textbf{Point 5:} Issue for presentation. It would be clearer if the authors could arrange the figures according to the sequence of the their occurrence in text. 
% For example, Figure 4 is mentioned first, yet it comes after Figure 2 and 3. This would cause confusion to readers.


\noindent \textbf{Reply:} Thanks a lot for the helpful suggestion! We have already made arrangements for better presentation in our revised manuscript.



\section{Response to Reviewer \#2}
\label{sec:reviewer_2}

\noindent \textbf{Point 1:} The motivation is not well summarized and explained. Why is this method better than CoOp in a high-level perspective? When compared with CoOp in the abstract and introduction, the motivation from the author is ``another alternative path'', without in-depth or intuitive explanation. 
% This problem may make the novelty of the work less appealing.

\noindent \textbf{Reply:} Our motivation is based on the limitations of the CoOp method, such as the increased training costs due to the need to compute and store gradients for the entire encoder during back-propagation. Our aim was to propose a method that is more efficient and also enhances downstream visual recognition performances. Specifically, our residual method requires less computation and storage during training compared to CoOp, resulting in faster and more efficient training. Additionally, our approach achieves state-of-the-art performance on several benchmarks, indicating that it can improve the zero-shot/few-shot classification performance of vision-language models. Overall, our method offers improved performance and efficiency compared to CoOp.

\vspace{5pt}
\noindent \textbf{Point 2:} The related works of this paper need to be improved. Some previous works are misunderstood. For example, CPT uses natural language prompts instead of optimizing continuous prompts. In fact, natural language prompts for vision-language pretrained models are the earliest attempts and are still in rapid development\cite{tsimpoukelli2021multimodal,alayrac2022flamingo,wang2022simvlm,yao-etal-2022-pevl}. It will be better to include more related works in this direction.

\noindent \textbf{Reply:} Thanks for pointing out this! We have already corrected the description for the CPT approach. In addition, we have improved the related work section of our paper to include more suggested related works in the field of ``natural language prompts for vision-language pretrained models''.


\section{Response to Reviewer \#3}
\label{sec:reviewer_3}

\noindent \textbf{Point 1:} The current version of CLIP-Adapter is only implemented in the vision backbone, as it found adapting the vision module is more effective than adapting the language module. From my perspective, the effectiveness of ``CLIP-Adapter'' comes from tuning the vision backbone instead of designing a new ``Adapter'' module. Therefore, to prove the effectiveness of CLIP-Adapter, it will appreciate that paper should investigate different vision backbone adaptation strategies as a strong baseline.  (1) inserting the original Adapter in the vision backbone. (2) ``Language-Init Adaptation''~\cite{li2022elevater} should be a better linear probing baseline rather than ``CLIP linear probing''.

\noindent \textbf{Reply:} (1) As shown in Table 3, we conducted the experiment for evaluating the effectiveness of CLIP-Adapter compared to the original Adapter (houlsby-adapter~\cite{houlsby2019parameter}) in the vision backbone. We believe that this comparison provides a stronger baseline for assessing the effectiveness and efficiency of our proposed method and proves our design choice of CLIP-Adapter.

\noindent (2) To ensure consistency and comparability with the results presented in the CoOp paper, we adopted the linear-probe evaluation approach outlined in their work. However, we acknowledge the suggestion to include a comparison with Language-Init Adaptation as a stronger baseline for our proposed CLIP-Adapter method.  Therefore, we add this comparison in Table 4 of our paper to further demonstrate the effectiveness of our approach. From the table, we can see that compared to training additional pro-jection layers initialized by randomness on the text
encoder, our CLIP-Adapter achieves higher classification accuracy, since the residual design can largely preserve the pretrained knowledge in CLIP.


%%===========================================================================================%%
%% If you are submitting to one of the Nature Portfolio journals, using the eJP submission   %%
%% system, please include the references within the manuscript file itself. You may do this  %%
%% by copying the reference list from your .bbl file, paste it into the main manuscript .tex %%
%% file, and delete the associated \verb+\bibliography+ commands.                            %%
%%===========================================================================================%%
\bibliographystyle{sn-basic}
\bibliography{custom}% common bib file
%% if required, the content of .bbl file can be included here once bbl is generated
%%\input sn-article.bbl

%% Default %%
%%\input sn-sample-bib.tex%

\end{document}
